\section{Dense endpoint progressions certified by two overlaps}
\label{sec:lcs-endpoint}

The 65-marker limitation of Section~\ref{sec:lcs-sparse} suggests retaining
an entire arithmetic progression of candidate lengths. This removes both
the fixed marker-count limit and the short-period limit for a useful family,
while preserving the general exact occurrence matcher as fallback.

\subsection{Certifying an occurrence grid without enumeration}
For a signed letter $a$, summarize its positions in a grammar word by
$(f,l,d,c)$: first position, last position, gcd of successive gaps, and exact
cardinality. Empty sets have $c=0$; singletons have $d=0$. At a concatenation
$UV$, shift the right positions by $|U|$. If both sets are nonempty, the new
gap gcd is
\[
 \gcd(d_U,d_V,|U|+f_V-l_U).
\]
The first and last positions and count combine directly. Every position is
congruent to $f$ modulo $d$. Since positions are distinct, the set fills its
entire grid if and only if
\[
 l-f=d(c-1).
\]
This criterion also covers a singleton. Absence, a certified progression,
and a nonprogression are distinct outcomes. The traversal uses the reachable
grammar, including binary arithmetic on exponentially large counts and gaps;
it never expands those counts into lists.

Every suffix/prefix overlap of $X,Y$ has a length in both necessary sets
\[
 \{|X|-i:X_i=\operatorname{first}(Y)\},\qquad
 \{j+1:Y_j=\operatorname{last}(X)\}.
\]
If both sets are progressions, intersect them using exact CRT arithmetic.
If only one is a progression, it still supplies a necessary candidate set.
Clip to $1\le k\le\min(|X|,|Y|)$. An absent endpoint or empty intersection
proves that there are no overlaps. If neither set is a progression, decline
this stage. No conclusion is drawn from failure to find a grid.

\subsection{The two-largest-overlaps certificate}
Suppose the resulting candidate progression contains at least two lengths,
with maximum $L$ and step $d>0$. Check exactly, using compressed equality,
that both $L$ and $L-d$ are overlaps. If so, the common word
\[
 W=\operatorname{suffix}_L(X)=\operatorname{prefix}_L(Y)
\]
has a border of length $L-d$, hence period $d$. Repeated application of
that period implies that every positive $L-jd$ is also a border of $W$.
It follows that every member of the candidate progression is an overlap.
Conversely, endpoint necessity already ensures that every overlap belongs
to that progression. Thus two exact equalities certify the entire answer.
The progression need not begin at $d$, and its cardinality need not fit a
machine integer or any fixed enumeration allowance.

For a singleton, one equality decides it. If either of the two checks for
a nonsingleton fails, the stage declines: smaller candidates may still be
valid. The complete matcher then runs. This distinction is essential;
a failed largest overlap is not a certificate of an empty overlap set.

For $s$ reachable rules, grammar height $H$, and maximum integer bit length
$b$, the two grid passes cost $O(s\log s)$ word operations for reachability
and a polynomial amount of $b$-bit arithmetic per rule. The CRT intersection
has polynomial bit cost. At most four slices and two equality queries remain,
adding $O(H)$ grammar nodes. If $E(t,b)$ bounds the deterministic equality
primitive, the equality term is at most $2E(s+O(H),b)$. These costs remain
charged to the common work and node allowances; a limit is inconclusive.
The implementation tries this route after the cheaper uniform, short-period,
and sparse-endpoint routes, for operands of length at least 128.

Set $B=(ab)^Nc$ and $X=Y=B^q$. The occurrences of $c$ give precisely the
length progression $(|B|,|B|,q)$. For $q\ge2$, the two largest candidates
are true overlaps, so the complete answer is certified without enumerating
$q$. In particular both $N$ and $q$ may equal $2^{500}$. This is a local
compressed-string theorem. It does not bound the number of relator moves,
intermediate grammar growth, or total unknot-recognition time.

\subsection{Validation and measurements}
The four new tests compare occurrence grids with literal sets on all binary
words of length at most seven and 300 random signed words; compare complete
overlap sets and critical-extension maxima on 40 randomized repeated blocks;
exercise disjoint grids and both failed-check fallbacks; and certify the
$N=q=2^{500}$ family under a 100000-work allowance with expansion and generic
occurrence tables forbidden. Existing generic-matcher coverage explicitly
bypasses the independent fast stages. The full suite passes all 701 tests
with Regina 7.4.1 in 166.342 seconds.

\begin{center}\small
\begin{tabular}{@{}lrrr@{}}
\toprule
Scope & Before & A/A & Grid \\
\midrule
Small 14, summed medians & 197.850 & 198.286 & 198.954 \\
Gordian, full query & 2278.670 & 1997.217 & 3209.060 \\
\bottomrule
\end{tabular}
\end{center}
Median milliseconds; all full knot queries complete.

\begin{center}\small
\begin{tabular}{@{}llrrr@{}}
\toprule
Family & $h$ & Before & A/A & Grid \\
\midrule
Prefix & 8 & 0.103 & 0.104 & 0.105 \\
Prefix & 32 & 0.272 & 0.271 & 0.271 \\
Prefix & 500 & 3.714 & 3.708 & 3.664 \\
Interior & 8 & 0.181 & 0.158 & 0.158 \\
Interior & 32 & 0.434 & 0.405 & 0.431 \\
Interior & 500 & 5.345 & 5.375 & 5.359 \\
Equal pairs & 8 & 2.301 & 2.303 & 2.311 \\
Equal pairs & 32 & 2.608 & 2.635 & 2.612 \\
Equal pairs & 500 & 8.469 & 8.472 & 8.490 \\
Full overlap & 8 & 0.097 & 0.095 & 0.094 \\
Full overlap & 32 & 0.226 & 0.200 & 0.199 \\
Full overlap & 500 & 2.506 & 2.467 & 2.485 \\
Shifted 32 & 8 & 0.459 & 0.441 & 0.440 \\
Shifted 32 & 32 & 0.619 & 0.621 & 0.633 \\
Shifted 32 & 500 & 4.279 & 4.310 & 4.319 \\
Late defect & 8 & 0.119 & 0.119 & 0.120 \\
Late defect & 32 & 0.250 & 0.275 & 0.252 \\
Late defect & 500 & 3.076 & 3.097 & 3.105 \\
Markers 17 & 8 & 0.871 & 0.841 & 0.833 \\
Markers 17 & 32 & 1.059 & 1.033 & 1.066 \\
Markers 17 & 500 & 3.759 & 3.871 & 3.877 \\
Markers 64 & 8 & 11.050 & 11.055 & 11.092 \\
Markers 64 & 32 & 11.775 & 11.798 & 11.789 \\
Markers 64 & 500 & 14.938 & 15.058 & 15.126 \\
Markers 65 & 8 & 213.510 & 209.198 & 0.675 \\
Markers 65 & 32 & limit & limit & 0.942 \\
Markers 65 & 500 & limit & limit & 4.810 \\
Markers power & 8 & 314.478 & 321.601 & 0.917 \\
Markers power & 32 & limit & limit & 3.399 \\
Markers power & 500 & limit & limit & 52.894 \\
Broken grid & 8 & 255.670 & 271.733 & 259.146 \\
Broken grid & 32 & limit & limit & limit \\
Broken grid & 500 & limit & limit & limit \\
\bottomrule
\end{tabular}
\end{center}
Median milliseconds including construction; limits are censored. Full overlap enumerates all lengths in compressed progressions, not all occurrences.

\begin{center}\small
\begin{tabular}{@{}llrrr@{}}
\toprule
Family & $h$ & Before & A/A & Grid \\
\midrule
Partial & 8 & 0.311 & 0.325 & 0.314 \\
Partial & 32 & 0.665 & 0.650 & 0.662 \\
Partial & 500 & 7.000 & 7.059 & 7.310 \\
Equal-pair cyclic & 8 & 0.510 & 0.505 & 0.508 \\
Equal-pair cyclic & 32 & 0.835 & 0.837 & 0.842 \\
Equal-pair cyclic & 500 & 7.158 & 7.248 & 7.164 \\
Quotient & 8 & 0.336 & 0.335 & 0.339 \\
Quotient & 32 & 0.979 & 0.978 & 0.963 \\
Quotient & 500 & 13.937 & 13.636 & 13.726 \\
\bottomrule
\end{tabular}
\end{center}
Median milliseconds including construction; limits are censored. Local operations do not produce knot verdicts.


For $q=65$ and $N=2^8$, median time falls from 213.510 milliseconds
(209.198 A/A) to 0.675; charged work falls from 259320 to 860 and arena
nodes from 425 to 26. At $N=2^{500}$ both old arms exhaust two million
work units; the new query completes in 4.810 milliseconds, with 518 nodes,
9224 work units, two exact equalities and no occurrence table.
For $q=N=2^{500}$ the complete answer is one progression containing
$2^{500}$ lengths. Its median is 52.894 milliseconds, using 2003 nodes,
48292 work units and two equalities. Both controls are limited.
For $q=N=2^8$, completed medians are 314.478, 321.601 and 0.917 milliseconds.

The negative control changes the first block to $(ba)^Nc$ while retaining
64 ordinary blocks, and compares that word with itself. Its endpoint grid
is regular, but the second-largest proposed overlap fails. Only the full
word is an overlap. At $N=2^8$, the fallback completes in 255.670 milliseconds
before, 271.733 A/A and 259.146 after. The unsuccessful certificate adds
210 work units and six nodes. At $N=2^{32}$ and $N=2^{500}$ every arm is
still limited. A regular endpoint grid alone is therefore insufficient.

All 225 whole-knot queries preserve their certificate digests and none
activates this stage. Small-fourteen sums of medians are 197.850, 198.286
and 198.954 milliseconds. Gordian medians are 2.279, 1.997 and 3.209 seconds;
this batch is noisy and does not demonstrate a recognition speedup.
The 135 local relator operations retain identical traces; the $N=2^{500}$
quotient medians are 13.937, 13.636 and 13.726 milliseconds. Earlier phase
and sparse routes likewise retain their counters on the corresponding controls.

The archived baseline is \code{8170a64c7}. A/A repeats its matcher inside
the common current arena and recognition/replay pipeline. Seed 3015 shuffles
three arms in each of five measured rounds, after one excluded warmup.
The 495 string queries include construction and local answer validation;
no progression with exponentially many entries is enumerated for validation.
Local allowances are two million work units and 100000 arena/table cells.
Whole recognition allows 18 seconds, with a 15-second, twenty-million-work
group stage. Limited samples remain censored, with no completed-time ratio.
Raw measurements are in \path{results/lcs_endpoint_progressions_20261008.json};
the test and benchmark logs are under \path{synthesis/data/lcs-endpoint-*}.
